\section{Properties of the spectral family of a bounded self-adjoint linear operator}

1 Again, let us try to deduce the finite-dim case 
2 The compact case 
3 The multiplication operator case 
4
5
6 The general multiplication operator
7 More investigations of the spectrum
8
9
10

\begin{exercise}{1}
What can we conclude from Theorem 9.11-1 in the case of a Hermitian matrix?
\end{exercise}
\begin{proof}
fill
\end{proof} 

\begin{exercise}{2}
If $T$ in Theorem 9.11-1 is compact and has infinitely many eigenvalues, what can we conclude about $(E_k)$ from Theorems 9.11-1 and 9.11-2?
\end{exercise}
\begin{proof}
fill
\end{proof} 

\begin{exercise}{3}
Verify that the spectral family in exercise 9.9.7 satisfies the three theorems in the present section.
\end{exercise}
\begin{proof}
fill
\end{proof} 

\begin{exercise}{4}
We know that if $m$ in Theorem 9.2-1 is positive, then $T$ is positive.
How does this follow from the spectral representation $T = \int_{m-0}^M \lambda dE_\lambda$ in section 9.9?
\end{exercise}
\begin{proof}
fill
\end{proof} 

\begin{exercise}{5}
We know that the spectrum of a bounded self-adjoint linear operator is closed.
How does this follow from theorems in this section?
\end{exercise}
\begin{proof}
fill
\end{proof} 

\begin{exercise}{6}
Let $T: \ell^2 \to \ell^2$ be defined by $y = (y_n) = Tx$, where $x = (x_n)$, $y_n = \alpha_n x_n$ and $(\alpha_n)$ is any real sequence in a finite interval $[a, b]$.
Show that the corresponding spectral family $(E_\lambda)$ is defined by $\brackets{E_\lambda x, y} = \sum_{\alpha_j \leq \lambda} x_j \bar{y}_j$.
\end{exercise}
\begin{proof}
fill
\end{proof} 

\begin{exercise}{7 (Pure point spectrum)}
A bounded self-adjoint linear operator $T: H \to H$ on a Hilbert space $H \neq \set{0}$ is said to have a pure point spectrum or purely discrete spectrum if $T$ has an orthonormal set of eigenvectors which is total in $H$.
Illustrate with an example that this does not imply $\sigma_c(T)$ (so that this terminology, which is generally used, may confuse the beginner for a moment).
\end{exercise}
\begin{proof}
fill
\end{proof} 

\begin{exercise}{8}
Give examples of compact self-adjoint linear operators $T: \ell^2 \to \ell^2$ having a pure point spectrum such that the set of the nonzero eigenvalues 
(a) is a finite point set,
(b) is an infinite point set and the corresponding eigenvectors form a dense set in $\ell^2$,
(c) is an infinite point set and the corresponding eigenvectors span a subspace of $\ell^2$ such that the orthogonal complement of the closure of that subspace is finite dimensional,
(d) as in (c) but that complement is infinite dimensional.
In each case find a total orthonormal set of eigenvectors.
\end{exercise}
\begin{proof}
fill
\end{proof} 

\begin{exercise}{9 (Purely continuous spectrum)}
A bounded self-adjoint linear operator $T: H \to H$ on a Hilbert space $H \neq \set{0}$ is said to have a purely continuous spectrum if $T$ has no eigenvalues.
If $T$ is any bounded self-adjoint operator on $H$, show that there is a closed subspace $Y \subseteq H$ which reduces $T$ (see exercise 9.6.10) and is such that $T_1 = T|_Y$ has a pure point spectrum, whereas $T_2 = T|_Z, Z = ^\perp$ has a purely continuous spectrum.
(This reduction facilitated the investigation of $T$).
\end{exercise}
\begin{proof}
fill
\end{proof} 

\begin{exercise}{10}
What can we say about the spectral families $(E_{\lambda 1})$ and $(E_{\lambda 2})$ of $T_1$ and $T_2$ in exercise 9 in terms of the spectral family $(E_\lambda)$ of $T$?
\end{exercise}
\begin{proof}
fill
\end{proof} 
